% Define colors matching your preferred tcolorbox style
\definecolor{propbox}{RGB}{240, 248, 255}    % Light blue for definition context
\definecolor{propborder}{RGB}{0, 102, 204}   % Dark blue for border
\definecolor{setbox}{RGB}{204, 229, 255}    % Slightly darker blue for the set itself
\definecolor{pointcolor}{RGB}{0, 0, 0}        % Black for points

% Set common node style matching tcolorbox look
\tikzset{
    databox/.style={
        draw=propborder,
        fill=propbox,
        thick,
        rounded corners=2pt,
        align=center,
        inner sep=6pt,
        font=\small
    },
    point/.style={
        fill=pointcolor,
        circle,
        inner sep=1.8pt
    }
}

\begin{figure}[h]
\centering
\begin{tikzpicture}[
    >=Stealth,
    x=1.8cm, 
    thick,
    every node/.style={font=\small}
]

    % =====================================================
    % 1. TOP LINE: INFIMUM AND LOWER BOUNDS
    % =====================================================
    \draw[->, gray!60] (-4.5, 2) -- (1.5, 2) node[right, black] {$\mathbb{R}$};
    
    % The Set A (visualized as half-infinite for this context)
    \draw[line width=3.5pt, propborder!30] (-2.5, 2) -- (1.0, 2);
    \node[anchor=south] at (-0.75, 2.1) {$\textbf{Set } A$};

    % Lower Bounds Region
    \draw[line width=1.5pt, propborder] (-4.3, 3.0) -- (-2.5, 3.0);
    \fill (-2.5, 3.0) circle (1.5pt); 
    \node[anchor=east] at (-4.3, 3.0) {$\mathbb{L}(A)$ (Lower Bounds)};

    % Infimum Annotation
    \draw[<-, thin, gray] (-2.5, 2.1) -- (-2.5, 3.5) node[above] {$\inf(A)$};
    \node[anchor=north east, scale=0.8, align=right] at (-2.6, 2.9) {The Greatest \\ Lower Bound};

    % Minimum condition
    \draw[|-|, thick, red] (-2.5, 1.4) -- (-2.5, 1.8) node[midway, left=0.1cm, black] {$\min(A)$};
    \node[anchor=north west, scale=0.8, align=left] at (-2.4, 1.7) {Exists iff \\ $\inf(A) \in A$};

    % =====================================================
    % 2. BOTTOM LINE: SUPREMUM AND UPPER BOUNDS
    % =====================================================
    \begin{scope}[yshift=-2.5cm]
        \draw[->, gray!60] (-1.5, 2) -- (4.5, 2) node[right, black] {$\mathbb{R}$};
        
        % The Set A (visualized as half-infinite for this context)
        \draw[line width=3.5pt, propborder!30] (-1.0, 2) -- (2.5, 2);
        \node[anchor=south] at (0.75, 2.1) {$\textbf{Set } A$};

        % Upper Bounds Region
        \draw[line width=1.5pt, propborder] (2.5, 3.0) -- (4.3, 3.0);
        \fill (2.5, 3.0) circle (1.5pt);
        \node[anchor=west] at (4.3, 3.0) {$\mathbb{U}(A)$ (Upper Bounds)};

        % Supremum Annotation
        \draw[<-, thin, gray] (2.5, 2.1) -- (2.5, 3.5) node[above] {$\sup(A)$};
        \node[anchor=north west, scale=0.8, align=left] at (2.6, 2.9) {The Least \\ Upper Bound};

        % Maximum condition
        \draw[|-|, thick, red] (2.5, 1.4) -- (2.5, 1.8) node[midway, right=0.1cm, black] {$\max(A)$};
        \node[anchor=north east, scale=0.8, align=right] at (2.4, 1.7) {Exists iff \\ $\sup(A) \in A$};
    \end{scope}

\end{tikzpicture}
\caption{Separated geometric relationships for lower and upper bounds.}
\end{figure}




\begin{figure}[h]
\centering
\begin{tikzpicture}[
    >=Stealth,
    x=4.5cm,
    y=1.0cm,
    thick,
    every node/.style={font=\small}
]
    % =====================================================
    % VERTICAL LANE MAP
    %   y = +1.4  : "Set A" label (high, above line)
    %   y = +0.7  : "s = sup(A)" label (mid, above line)
    %   y = +0.0  : The real line itself
    %   y = -0.5  : "a ∈ A" label (just below line)
    %   y = -1.2  : epsilon probe arrow
    %   y = -1.7  : "s - ε" label
    %   y = -2.6  : "Guaranteed Point / a ∈ (s-ε, s]" label
    %   y = -3.4  : brace
    %   y = -3.9  : brace caption
    % =====================================================

    % --------------------------------------------------
    % 1. REAL LINE
    % --------------------------------------------------
    \draw[->, gray!60] (-1.3, 0) -- (1.5, 0)
        node[right, black] {$\mathbb{R}$};

    % --------------------------------------------------
    % 2. SET A  (thick cyan bar, x ∈ [-1.2, 0])
    % --------------------------------------------------
    \draw[line width=5pt, cyan!40] (-1.2, 0) -- (0, 0);

    % "Set A" label: high lane, centered over the bar
    \node[anchor=south] at (-0.6, 1.2)
        {\textbf{Set} $A$};
    % Thin leader line from label down to the bar
    \draw[thin, gray!70] (-0.6, 1.15) -- (-0.6, 0.08);

    % --------------------------------------------------
    % 3. SUPREMUM DOT AND LABEL
    % --------------------------------------------------
    \fill[black] (0, 0) circle (2pt);

    % "s = sup(A)" label: mid lane, directly above the dot
    \node[anchor=south] at (0, 0.55)
        {$s = \sup(A)$};
    % Short leader from label to dot
    \draw[thin, gray!70] (0, 0.52) -- (0, 0.08);

    % Full dashed vertical spine (reference column)
    \draw[thin, dashed, gray] (0, 0.5) -- (0, -3.5);

    % --------------------------------------------------
    % 4. UPPER BOUNDS REGION
    % --------------------------------------------------
    \draw[line width=2.5pt, blue!70] (0, 0) -- (1.4, 0);
    % Label sits above the bar, left-anchored from x=0.08
    \node[anchor=south west] at (0.08, 0.08)
        {$\mathbb{U}(A)$ \small(Upper Bounds)};

    % --------------------------------------------------
    % 5. WITNESS ELEMENT 'a'
    % --------------------------------------------------
    \coordinate (witness) at (-0.3, 0);
    \fill[black] (witness) circle (2pt);

    % "a ∈ A" label: dedicated below-line lane, clear of sup label
    \node[anchor=north] at (-0.3, -0.35)
        {$a \in A$};

    % --------------------------------------------------
    % 6. EPSILON PROBE  —  Lane y = -1.2
    % --------------------------------------------------
    % Dashed verticals dropping from the line to the probe
    \draw[thin, dashed, red] (-0.7,  0) -- (-0.7, -1.2);
    \draw[thin, dashed, red] ( 0,    0) -- ( 0,   -1.2);

    % The double-headed probe arrow
    \draw[<->, red] (-0.7, -1.2) -- (0, -1.2)
        node[midway, above=2pt, fill=white, inner sep=1pt]
        {\textcolor{red}{$\varepsilon$}};

    % s - ε dot and label — Lane y = -1.7
    \fill[red] (-0.7, 0) circle (1.5pt);
    \node[anchor=north] at (-0.7, -1.25)
        {$s - \varepsilon$};

    % --------------------------------------------------
    % 7. GUARANTEED POINT LABEL + ARROW  —  Lane y = -2.6
    % --------------------------------------------------
    % Arrow rises from the label up toward the witness dot
    % (stops 0.15 above the line so it doesn't pierce the dot)
    \draw[->, thin, blue] (-0.3, -2.45) -- (-0.3, 0.15);

    \node[blue, anchor=north, align=center] at (-0.3, -2.5)
        {\textbf{Guaranteed Point:}\\[3pt]
         $a \in (s-\varepsilon,\, s]$};

    % --------------------------------------------------
    % 8. BRACE + CAPTION  —  Lanes y = -3.4 / -3.9
    % --------------------------------------------------
    \draw[decorate,
          decoration={brace, amplitude=5pt, mirror, raise=4pt}]
        (-0.7, -3.35) -- (0, -3.35)
        node[midway, below=12pt, align=center, scale=0.85]
        {The neighborhood where\\points of $A$ \emph{must} exist};

\end{tikzpicture}
\caption{Refined visualization of the $\varepsilon$-characterization.
Labels occupy staggered vertical lanes with leader lines to eliminate
all horizontal collision, while the dashed spine preserves the geometric
correspondence between the real line and the probe below it.}
\end{figure}